% =========================================================
% ppt/src/minggu-02.tex
% Slide pertemuan 2
% Bab IV — Arsitektur Mikroprosesor 8086
%
% Sub-CPMK 1.1 | CPMK01 | CPL02
% =========================================================

\documentclass[aspectratio=169,11pt]{beamer}
% =========================================================
% ppt/src/preamble-ppt.tex
% Tema, warna, paket, dan makro bersama seluruh deck slide.
%
% Dipakai oleh setiap berkas minggu-NN.tex dengan pola:
%   \documentclass[aspectratio=169,11pt]{beamer}
%   % =========================================================
% ppt/src/preamble-ppt.tex
% Tema, warna, paket, dan makro bersama seluruh deck slide.
%
% Dipakai oleh setiap berkas minggu-NN.tex dengan pola:
%   \documentclass[aspectratio=169,11pt]{beamer}
%   % =========================================================
% ppt/src/preamble-ppt.tex
% Tema, warna, paket, dan makro bersama seluruh deck slide.
%
% Dipakai oleh setiap berkas minggu-NN.tex dengan pola:
%   \documentclass[aspectratio=169,11pt]{beamer}
%   \input{preamble-ppt}
%
% Palet warna disalin dari buku_ajar/config/styles.tex agar
% identitas visual slide sama dengan buku ajar.
% =========================================================

% ---------------------------------------------------------
% 1. Paket
% ---------------------------------------------------------
\usepackage[T1]{fontenc}
\usepackage[utf8]{inputenc}
\usepackage[indonesian]{babel}
\usepackage{lmodern}
\usepackage{tikz}
\usetikzlibrary{arrows.meta, positioning, calc, shapes.geometric}
\usepackage{booktabs}
\usepackage{array}
\usepackage{amsmath}
\usepackage{amsfonts}
\usepackage{pifont}          % \ding untuk tanda checklist
\usepackage{listings}

% ---------------------------------------------------------
% 2. Palet warna OBE (identik dengan buku)
% ---------------------------------------------------------
\definecolor{obeBlue}{HTML}{1F3864}
\definecolor{obeBlueLight}{HTML}{DCE4F5}
\definecolor{obeTeal}{HTML}{0F6E6E}
\definecolor{obeTealLight}{HTML}{D9F0EE}
\definecolor{obeGreen}{HTML}{2E6B34}
\definecolor{obeGreenLight}{HTML}{DEEFDF}
\definecolor{obeAmber}{HTML}{8A5A00}
\definecolor{obeAmberLight}{HTML}{FDF1D6}
\definecolor{obeMaroon}{HTML}{7A2233}
\definecolor{obeMaroonLight}{HTML}{F7E1E5}
\definecolor{obeGray}{HTML}{4A4A4A}
\definecolor{obeGrayLight}{HTML}{F0F0F0}
\definecolor{codeBack}{HTML}{F7F7F4}
\definecolor{codeFrame}{HTML}{C9C9C0}

% ---------------------------------------------------------
% 3. Identitas mata kuliah (dipakai seluruh deck)
% ---------------------------------------------------------
\newcommand{\NamaMK}{Bahasa Pemrograman Tingkat Rendah}
\newcommand{\NamaMKShort}{Bhs.\ Pemrograman Tingkat Rendah}
\newcommand{\KodeMK}{MWCA010}
\newcommand{\Prodi}{Program Studi Teknik Informatika}
\newcommand{\Fakultas}{Fakultas Teknologi Informasi}
\newcommand{\PerguruanTinggi}{Universitas Bale Bandung}
\newcommand{\DosenPengampu}{Cecep Suwanda, S.Si., M.Kom.}
\newcommand{\SemesterTeks}{Semester Ganjil}

% \PertemuanTeks dan \JudulBab ditetapkan di masing-masing
% berkas minggu-NN.tex sebelum \begin{document}.
\providecommand{\PertemuanTeks}{Pertemuan}

% ---------------------------------------------------------
% 4. Tema beamer
% ---------------------------------------------------------
\usetheme{default}
\setbeamertemplate{navigation symbols}{}

\setbeamercolor{normal text}{fg=black!88}
\setbeamercolor{structure}{fg=obeBlue}
\setbeamercolor{alerted text}{fg=obeMaroon}

% --- judul slide: pita biru penuh lebar ---
\setbeamercolor{frametitle}{fg=white,bg=obeBlue}
\setbeamerfont{frametitle}{series=\bfseries,size=\large}
\setbeamertemplate{frametitle}{%
    \nointerlineskip%
    \begin{beamercolorbox}[%
        wd=\paperwidth,ht=2.9ex,dp=1.5ex,
        leftskip=0.4cm,rightskip=0.4cm]{frametitle}%
        \usebeamerfont{frametitle}\insertframetitle%
    \end{beamercolorbox}%
    \nointerlineskip%
    \hbox{\textcolor{obeTeal}{\rule{\paperwidth}{1.1pt}}}%
}

% --- kaki slide: identitas mata kuliah, pertemuan, nomor ---
\setbeamercolor{footline}{fg=white,bg=obeGray}
\setbeamerfont{footline}{size=\scriptsize,family=\sffamily}
\setbeamertemplate{footline}{%
    \leavevmode%
    \hbox{%
        \begin{beamercolorbox}[wd=0.42\paperwidth,ht=2.1ex,dp=1.1ex,
            leftskip=0.4cm]{footline}%
            \usebeamerfont{footline}%
            \KodeMK\ \textbullet\ \NamaMKShort%
        \end{beamercolorbox}%
        \begin{beamercolorbox}[wd=0.32\paperwidth,ht=2.1ex,dp=1.1ex,
            center]{footline}%
            \usebeamerfont{footline}\PertemuanTeks%
        \end{beamercolorbox}%
        \begin{beamercolorbox}[wd=0.26\paperwidth,ht=2.1ex,dp=1.1ex,
            rightskip=0.4cm]{footline}%
            \usebeamerfont{footline}%
            \insertframenumber\,/\,\inserttotalframenumber%
        \end{beamercolorbox}%
    }%
    \vskip0pt%
}

% --- daftar ---
\setbeamercolor{itemize item}{fg=obeTeal}
\setbeamercolor{itemize subitem}{fg=obeTeal}
\setbeamercolor{enumerate item}{fg=obeBlue}
\setbeamertemplate{itemize item}{\raise0.2ex\hbox{\tiny$\blacksquare$}}
\setbeamertemplate{itemize subitem}{\raise0.2ex\hbox{\scriptsize$\bullet$}}

% --- kotak (pengganti tcolorbox pada buku) ---
\setbeamercolor{block title}{bg=obeBlueLight,fg=obeBlue}
\setbeamercolor{block body}{bg=obeBlueLight!45,fg=black!88}
\setbeamercolor{block title alerted}{bg=obeMaroonLight,fg=obeMaroon}
\setbeamercolor{block body alerted}{bg=obeMaroonLight!45,fg=black!88}
\setbeamercolor{block title example}{bg=obeTealLight,fg=obeTeal}
\setbeamercolor{block body example}{bg=obeTealLight!45,fg=black!88}

% ---------------------------------------------------------
% 5. Makro notasi (sama persis dengan buku)
%
% Catatan: \note TIDAK didefinisikan di sini karena beamer
% sudah memakai \note untuk catatan pembicara. Untuk catatan
% di slide dipakai \catatan.
% ---------------------------------------------------------
\newcommand{\code}[1]{\texttt{#1}}
\newcommand{\asm}[1]{\texttt{#1}}
\newcommand{\reg}[1]{\texttt{#1}}
\newcommand{\flag}[1]{\texttt{#1}}
\newcommand{\dir}[1]{\texttt{#1}}
\newcommand{\file}[1]{\texttt{#1}}
\newcommand{\eng}[1]{\textit{#1}}

\newcommand{\hex}[1]{\texttt{#1h}}
\newcommand{\bin}[1]{\texttt{#1b}}
\newcommand{\dec}[1]{\texttt{#1}}

\newcommand{\important}[1]{\textbf{Penting:} #1}
\newcommand{\catatan}[1]{\textbf{Catatan:} #1}
\newcommand{\term}[1]{\textbf{#1}}

\newcommand{\CPL}[1]{\textbf{CPL#1}}
\newcommand{\CPMK}[1]{\textbf{CPMK#1}}
\newcommand{\SubCPMK}[1]{\textbf{Sub-CPMK #1}}

\newcommand{\checkbox}{\ding{113}\hspace{0.35em}}
\newcommand{\checked}{\ding{51}\hspace{0.35em}}
\newcommand{\crossed}{\ding{55}\hspace{0.35em}}

\newcommand{\flowarrow}{\ensuremath{\rightarrow}}

% ---------------------------------------------------------
% 6. Makro tata letak slide
% ---------------------------------------------------------

% Keterangan sumber materi (ketertelusuran ke buku)
\newcommand{\Sumber}[1]{%
    \vfill
    {\scriptsize\color{obeGray}Sumber: #1\par}%
}

% Frame judul deck: \JudulDeck{<judul bab>}{<pertemuan>}
\newcommand{\JudulDeck}[2]{%
\begin{frame}[plain]
    \vspace{0.25cm}
    {\color{obeBlue}\sffamily\bfseries\small \NamaMK\ (\KodeMK)\par}
    \vspace{0.12cm}
    {\color{obeGray}\sffamily\scriptsize
     \Prodi\ \textbullet\ \Fakultas\ \textbullet\ \PerguruanTinggi\par}
    \vfill
    {\color{obeBlue}\sffamily\LARGE\bfseries #1\par}
    \vspace{0.4cm}
    {\color{obeTeal}\sffamily\large #2\par}
    \vspace{0.45cm}
    {\color{obeTeal}\rule{\textwidth}{1pt}}
    \vfill
    {\color{obeGray}\sffamily\small Dosen Pengampu\par}
    \vspace{0.08cm}
    {\sffamily\bfseries \DosenPengampu\par}
    \vspace{0.2cm}
    {\color{obeGray}\sffamily\small \SemesterTeks\par}
    \vspace{0.25cm}
\end{frame}
}

% Frame pemisah bagian
\newcommand{\PemisahBagian}[1]{%
\begin{frame}[plain]
    \vfill
    \begin{center}
        {\color{obeTeal}\sffamily\small\bfseries BAGIAN\par}
        \vspace{0.35cm}
        {\color{obeBlue}\sffamily\Large\bfseries #1\par}
        \vspace{0.35cm}
        {\color{obeTeal}\rule{0.55\textwidth}{0.9pt}}
    \end{center}
    \vfill
\end{frame}
}

% ---------------------------------------------------------
% 7. Gaya penampil kode
% ---------------------------------------------------------
\input{listing-ppt}

%
% Palet warna disalin dari buku_ajar/config/styles.tex agar
% identitas visual slide sama dengan buku ajar.
% =========================================================

% ---------------------------------------------------------
% 1. Paket
% ---------------------------------------------------------
\usepackage[T1]{fontenc}
\usepackage[utf8]{inputenc}
\usepackage[indonesian]{babel}
\usepackage{lmodern}
\usepackage{tikz}
\usetikzlibrary{arrows.meta, positioning, calc, shapes.geometric}
\usepackage{booktabs}
\usepackage{array}
\usepackage{amsmath}
\usepackage{amsfonts}
\usepackage{pifont}          % \ding untuk tanda checklist
\usepackage{listings}

% ---------------------------------------------------------
% 2. Palet warna OBE (identik dengan buku)
% ---------------------------------------------------------
\definecolor{obeBlue}{HTML}{1F3864}
\definecolor{obeBlueLight}{HTML}{DCE4F5}
\definecolor{obeTeal}{HTML}{0F6E6E}
\definecolor{obeTealLight}{HTML}{D9F0EE}
\definecolor{obeGreen}{HTML}{2E6B34}
\definecolor{obeGreenLight}{HTML}{DEEFDF}
\definecolor{obeAmber}{HTML}{8A5A00}
\definecolor{obeAmberLight}{HTML}{FDF1D6}
\definecolor{obeMaroon}{HTML}{7A2233}
\definecolor{obeMaroonLight}{HTML}{F7E1E5}
\definecolor{obeGray}{HTML}{4A4A4A}
\definecolor{obeGrayLight}{HTML}{F0F0F0}
\definecolor{codeBack}{HTML}{F7F7F4}
\definecolor{codeFrame}{HTML}{C9C9C0}

% ---------------------------------------------------------
% 3. Identitas mata kuliah (dipakai seluruh deck)
% ---------------------------------------------------------
\newcommand{\NamaMK}{Bahasa Pemrograman Tingkat Rendah}
\newcommand{\NamaMKShort}{Bhs.\ Pemrograman Tingkat Rendah}
\newcommand{\KodeMK}{MWCA010}
\newcommand{\Prodi}{Program Studi Teknik Informatika}
\newcommand{\Fakultas}{Fakultas Teknologi Informasi}
\newcommand{\PerguruanTinggi}{Universitas Bale Bandung}
\newcommand{\DosenPengampu}{Cecep Suwanda, S.Si., M.Kom.}
\newcommand{\SemesterTeks}{Semester Ganjil}

% \PertemuanTeks dan \JudulBab ditetapkan di masing-masing
% berkas minggu-NN.tex sebelum \begin{document}.
\providecommand{\PertemuanTeks}{Pertemuan}

% ---------------------------------------------------------
% 4. Tema beamer
% ---------------------------------------------------------
\usetheme{default}
\setbeamertemplate{navigation symbols}{}

\setbeamercolor{normal text}{fg=black!88}
\setbeamercolor{structure}{fg=obeBlue}
\setbeamercolor{alerted text}{fg=obeMaroon}

% --- judul slide: pita biru penuh lebar ---
\setbeamercolor{frametitle}{fg=white,bg=obeBlue}
\setbeamerfont{frametitle}{series=\bfseries,size=\large}
\setbeamertemplate{frametitle}{%
    \nointerlineskip%
    \begin{beamercolorbox}[%
        wd=\paperwidth,ht=2.9ex,dp=1.5ex,
        leftskip=0.4cm,rightskip=0.4cm]{frametitle}%
        \usebeamerfont{frametitle}\insertframetitle%
    \end{beamercolorbox}%
    \nointerlineskip%
    \hbox{\textcolor{obeTeal}{\rule{\paperwidth}{1.1pt}}}%
}

% --- kaki slide: identitas mata kuliah, pertemuan, nomor ---
\setbeamercolor{footline}{fg=white,bg=obeGray}
\setbeamerfont{footline}{size=\scriptsize,family=\sffamily}
\setbeamertemplate{footline}{%
    \leavevmode%
    \hbox{%
        \begin{beamercolorbox}[wd=0.42\paperwidth,ht=2.1ex,dp=1.1ex,
            leftskip=0.4cm]{footline}%
            \usebeamerfont{footline}%
            \KodeMK\ \textbullet\ \NamaMKShort%
        \end{beamercolorbox}%
        \begin{beamercolorbox}[wd=0.32\paperwidth,ht=2.1ex,dp=1.1ex,
            center]{footline}%
            \usebeamerfont{footline}\PertemuanTeks%
        \end{beamercolorbox}%
        \begin{beamercolorbox}[wd=0.26\paperwidth,ht=2.1ex,dp=1.1ex,
            rightskip=0.4cm]{footline}%
            \usebeamerfont{footline}%
            \insertframenumber\,/\,\inserttotalframenumber%
        \end{beamercolorbox}%
    }%
    \vskip0pt%
}

% --- daftar ---
\setbeamercolor{itemize item}{fg=obeTeal}
\setbeamercolor{itemize subitem}{fg=obeTeal}
\setbeamercolor{enumerate item}{fg=obeBlue}
\setbeamertemplate{itemize item}{\raise0.2ex\hbox{\tiny$\blacksquare$}}
\setbeamertemplate{itemize subitem}{\raise0.2ex\hbox{\scriptsize$\bullet$}}

% --- kotak (pengganti tcolorbox pada buku) ---
\setbeamercolor{block title}{bg=obeBlueLight,fg=obeBlue}
\setbeamercolor{block body}{bg=obeBlueLight!45,fg=black!88}
\setbeamercolor{block title alerted}{bg=obeMaroonLight,fg=obeMaroon}
\setbeamercolor{block body alerted}{bg=obeMaroonLight!45,fg=black!88}
\setbeamercolor{block title example}{bg=obeTealLight,fg=obeTeal}
\setbeamercolor{block body example}{bg=obeTealLight!45,fg=black!88}

% ---------------------------------------------------------
% 5. Makro notasi (sama persis dengan buku)
%
% Catatan: \note TIDAK didefinisikan di sini karena beamer
% sudah memakai \note untuk catatan pembicara. Untuk catatan
% di slide dipakai \catatan.
% ---------------------------------------------------------
\newcommand{\code}[1]{\texttt{#1}}
\newcommand{\asm}[1]{\texttt{#1}}
\newcommand{\reg}[1]{\texttt{#1}}
\newcommand{\flag}[1]{\texttt{#1}}
\newcommand{\dir}[1]{\texttt{#1}}
\newcommand{\file}[1]{\texttt{#1}}
\newcommand{\eng}[1]{\textit{#1}}

\newcommand{\hex}[1]{\texttt{#1h}}
\newcommand{\bin}[1]{\texttt{#1b}}
\newcommand{\dec}[1]{\texttt{#1}}

\newcommand{\important}[1]{\textbf{Penting:} #1}
\newcommand{\catatan}[1]{\textbf{Catatan:} #1}
\newcommand{\term}[1]{\textbf{#1}}

\newcommand{\CPL}[1]{\textbf{CPL#1}}
\newcommand{\CPMK}[1]{\textbf{CPMK#1}}
\newcommand{\SubCPMK}[1]{\textbf{Sub-CPMK #1}}

\newcommand{\checkbox}{\ding{113}\hspace{0.35em}}
\newcommand{\checked}{\ding{51}\hspace{0.35em}}
\newcommand{\crossed}{\ding{55}\hspace{0.35em}}

\newcommand{\flowarrow}{\ensuremath{\rightarrow}}

% ---------------------------------------------------------
% 6. Makro tata letak slide
% ---------------------------------------------------------

% Keterangan sumber materi (ketertelusuran ke buku)
\newcommand{\Sumber}[1]{%
    \vfill
    {\scriptsize\color{obeGray}Sumber: #1\par}%
}

% Frame judul deck: \JudulDeck{<judul bab>}{<pertemuan>}
\newcommand{\JudulDeck}[2]{%
\begin{frame}[plain]
    \vspace{0.25cm}
    {\color{obeBlue}\sffamily\bfseries\small \NamaMK\ (\KodeMK)\par}
    \vspace{0.12cm}
    {\color{obeGray}\sffamily\scriptsize
     \Prodi\ \textbullet\ \Fakultas\ \textbullet\ \PerguruanTinggi\par}
    \vfill
    {\color{obeBlue}\sffamily\LARGE\bfseries #1\par}
    \vspace{0.4cm}
    {\color{obeTeal}\sffamily\large #2\par}
    \vspace{0.45cm}
    {\color{obeTeal}\rule{\textwidth}{1pt}}
    \vfill
    {\color{obeGray}\sffamily\small Dosen Pengampu\par}
    \vspace{0.08cm}
    {\sffamily\bfseries \DosenPengampu\par}
    \vspace{0.2cm}
    {\color{obeGray}\sffamily\small \SemesterTeks\par}
    \vspace{0.25cm}
\end{frame}
}

% Frame pemisah bagian
\newcommand{\PemisahBagian}[1]{%
\begin{frame}[plain]
    \vfill
    \begin{center}
        {\color{obeTeal}\sffamily\small\bfseries BAGIAN\par}
        \vspace{0.35cm}
        {\color{obeBlue}\sffamily\Large\bfseries #1\par}
        \vspace{0.35cm}
        {\color{obeTeal}\rule{0.55\textwidth}{0.9pt}}
    \end{center}
    \vfill
\end{frame}
}

% ---------------------------------------------------------
% 7. Gaya penampil kode
% ---------------------------------------------------------
% =========================================================
% ppt/src/listing-ppt.tex
% Gaya penampil kode untuk slide.
%
% Diturunkan dari buku_ajar/config/listings.tex, tetapi
% disesuaikan untuk slide: ukuran lebih kecil (karena frame
% 16:9 dan kode harus utuh terlihat) dan jarak lebih rapat.
% =========================================================

% ---------------------------------------------------------
% 1. Definisi bahasa MASM/TASM 8086 (salinan dari buku)
% ---------------------------------------------------------
\lstdefinelanguage{MASM}{
    morekeywords={
        MOV, XCHG, LEA, LDS, LES, PUSH, POP, PUSHA, POPA, PUSHF, POPF,
        LAHF, SAHF, XLAT, IN, OUT, CBW, CWD,
        ADD, ADC, SUB, SBB, INC, DEC, NEG, CMP, MUL, IMUL, DIV, IDIV,
        AAA, AAS, AAM, AAD, DAA, DAS,
        AND, OR, XOR, NOT, TEST, SHL, SHR, SAL, SAR, ROL, ROR, RCL, RCR,
        JMP, JE, JZ, JNE, JNZ, JG, JGE, JL, JLE, JA, JAE, JB, JBE,
        JC, JNC, JS, JNS, JO, JNO, JP, JPE, JNP, JPO, JCXZ,
        LOOP, LOOPE, LOOPZ, LOOPNE, LOOPNZ, CALL, RET, INT, IRET, INTO,
        HLT, NOP, CLI, STI, CLD, STD, CLC, STC, CMC, WAIT, ESC, LOCK,
        DB, DW, DD, DQ, DT, EQU, ORG, SEGMENT, ENDS, ASSUME, END,
        PROC, ENDP, MACRO, ENDM, LOCAL, PUBLIC, EXTRN, EXTERN,
        INCLUDE, INCLUDELIB, MODEL, STACK, CODE, DATA, CONST, ALIGN, EVEN,
        OFFSET, PTR, BYTE, WORD, DWORD, QWORD, TBYTE,
        NEAR, FAR, SHORT, LABEL, TYPE, SIZE, LENGTH, DUP, THIS,
        AX, BX, CX, DX, AL, AH, BL, BH, CL, CH, DL, DH,
        SP, BP, SI, DI, IP, CS, DS, SS, ES, FLAGS,
        CF, PF, AF, ZF, TF, IF, DF, OF
    },
    sensitive=false,
    morecomment=[l]{;},
    morestring=[b]',
    morestring=[b]"
}

% ---------------------------------------------------------
% 2. Gaya kode assembly untuk slide
% ---------------------------------------------------------
\lstdefinestyle{asmslide}{
    language=MASM,
    basicstyle=\ttfamily\scriptsize,
    keywordstyle=\color{obeBlue}\bfseries,
    commentstyle=\color{obeGreen}\itshape,
    stringstyle=\color{obeMaroon},
    numberstyle=\tiny\color{obeGray},
    numbers=left,
    numbersep=5pt,
    stepnumber=1,
    backgroundcolor=\color{codeBack},
    frame=single,
    rulecolor=\color{codeFrame},
    framesep=4pt,
    breaklines=true,
    breakatwhitespace=false,
    % CATATAN PENTING: opsi `postbreak` TIDAK boleh dipakai di sini.
    % Di beamer, `breaklines=true` bersama `postbreak=<kode>` membuat
    % listings gagal dirakit dengan galat
    %   "Missing number, treated as zero <to be read again> \kern".
    % Karena \lstset{style=asmslide} bersifat global, opsi postbreak
    % juga bocor ke gaya lain (mis. outslide/keluaran) dan merusaknya.
    % Di buku ajar (kelas `book`) justru sebaliknya: postbreak aman dan
    % dipakai. Jadi larangan ini khusus untuk slide.
    showstringspaces=false,
    tabsize=4,
    columns=fullflexible,
    keepspaces=true,
    upquote=true,
    xleftmargin=12pt,
    aboveskip=0.4em,
    belowskip=0.4em
}

% ---------------------------------------------------------
% 3. Gaya keluaran program / sesi debugger untuk slide
% ---------------------------------------------------------
\lstdefinestyle{outslide}{
    basicstyle=\ttfamily\scriptsize,
    keywordstyle=\color{black},
    commentstyle=\color{obeGray},
    stringstyle=\color{black},
    numbers=none,
    backgroundcolor=\color{obeGrayLight},
    frame=single,
    rulecolor=\color{codeFrame},
    framesep=4pt,
    breaklines=true,
    showstringspaces=false,
    columns=fullflexible,
    keepspaces=true,
    upquote=true,
    xleftmargin=4pt,
    aboveskip=0.4em,
    belowskip=0.4em
}

\lstset{style=asmslide}

% ---------------------------------------------------------
% 4. Environment ringkas untuk dipakai di dalam frame
%
% Frame yang memuat kedua environment ini WAJIB memakai
% opsi [fragile], karena isinya verbatim.
% ---------------------------------------------------------
\lstnewenvironment{kodeasm}[1][]{%
    \lstset{style=asmslide, #1}%
}{}

\lstnewenvironment{keluaran}[1][]{%
    \lstset{style=outslide, #1}%
}{}


%
% Palet warna disalin dari buku_ajar/config/styles.tex agar
% identitas visual slide sama dengan buku ajar.
% =========================================================

% ---------------------------------------------------------
% 1. Paket
% ---------------------------------------------------------
\usepackage[T1]{fontenc}
\usepackage[utf8]{inputenc}
\usepackage[indonesian]{babel}
\usepackage{lmodern}
\usepackage{tikz}
\usetikzlibrary{arrows.meta, positioning, calc, shapes.geometric}
\usepackage{booktabs}
\usepackage{array}
\usepackage{amsmath}
\usepackage{amsfonts}
\usepackage{pifont}          % \ding untuk tanda checklist
\usepackage{listings}

% ---------------------------------------------------------
% 2. Palet warna OBE (identik dengan buku)
% ---------------------------------------------------------
\definecolor{obeBlue}{HTML}{1F3864}
\definecolor{obeBlueLight}{HTML}{DCE4F5}
\definecolor{obeTeal}{HTML}{0F6E6E}
\definecolor{obeTealLight}{HTML}{D9F0EE}
\definecolor{obeGreen}{HTML}{2E6B34}
\definecolor{obeGreenLight}{HTML}{DEEFDF}
\definecolor{obeAmber}{HTML}{8A5A00}
\definecolor{obeAmberLight}{HTML}{FDF1D6}
\definecolor{obeMaroon}{HTML}{7A2233}
\definecolor{obeMaroonLight}{HTML}{F7E1E5}
\definecolor{obeGray}{HTML}{4A4A4A}
\definecolor{obeGrayLight}{HTML}{F0F0F0}
\definecolor{codeBack}{HTML}{F7F7F4}
\definecolor{codeFrame}{HTML}{C9C9C0}

% ---------------------------------------------------------
% 3. Identitas mata kuliah (dipakai seluruh deck)
% ---------------------------------------------------------
\newcommand{\NamaMK}{Bahasa Pemrograman Tingkat Rendah}
\newcommand{\NamaMKShort}{Bhs.\ Pemrograman Tingkat Rendah}
\newcommand{\KodeMK}{MWCA010}
\newcommand{\Prodi}{Program Studi Teknik Informatika}
\newcommand{\Fakultas}{Fakultas Teknologi Informasi}
\newcommand{\PerguruanTinggi}{Universitas Bale Bandung}
\newcommand{\DosenPengampu}{Cecep Suwanda, S.Si., M.Kom.}
\newcommand{\SemesterTeks}{Semester Ganjil}

% \PertemuanTeks dan \JudulBab ditetapkan di masing-masing
% berkas minggu-NN.tex sebelum \begin{document}.
\providecommand{\PertemuanTeks}{Pertemuan}

% ---------------------------------------------------------
% 4. Tema beamer
% ---------------------------------------------------------
\usetheme{default}
\setbeamertemplate{navigation symbols}{}

\setbeamercolor{normal text}{fg=black!88}
\setbeamercolor{structure}{fg=obeBlue}
\setbeamercolor{alerted text}{fg=obeMaroon}

% --- judul slide: pita biru penuh lebar ---
\setbeamercolor{frametitle}{fg=white,bg=obeBlue}
\setbeamerfont{frametitle}{series=\bfseries,size=\large}
\setbeamertemplate{frametitle}{%
    \nointerlineskip%
    \begin{beamercolorbox}[%
        wd=\paperwidth,ht=2.9ex,dp=1.5ex,
        leftskip=0.4cm,rightskip=0.4cm]{frametitle}%
        \usebeamerfont{frametitle}\insertframetitle%
    \end{beamercolorbox}%
    \nointerlineskip%
    \hbox{\textcolor{obeTeal}{\rule{\paperwidth}{1.1pt}}}%
}

% --- kaki slide: identitas mata kuliah, pertemuan, nomor ---
\setbeamercolor{footline}{fg=white,bg=obeGray}
\setbeamerfont{footline}{size=\scriptsize,family=\sffamily}
\setbeamertemplate{footline}{%
    \leavevmode%
    \hbox{%
        \begin{beamercolorbox}[wd=0.42\paperwidth,ht=2.1ex,dp=1.1ex,
            leftskip=0.4cm]{footline}%
            \usebeamerfont{footline}%
            \KodeMK\ \textbullet\ \NamaMKShort%
        \end{beamercolorbox}%
        \begin{beamercolorbox}[wd=0.32\paperwidth,ht=2.1ex,dp=1.1ex,
            center]{footline}%
            \usebeamerfont{footline}\PertemuanTeks%
        \end{beamercolorbox}%
        \begin{beamercolorbox}[wd=0.26\paperwidth,ht=2.1ex,dp=1.1ex,
            rightskip=0.4cm]{footline}%
            \usebeamerfont{footline}%
            \insertframenumber\,/\,\inserttotalframenumber%
        \end{beamercolorbox}%
    }%
    \vskip0pt%
}

% --- daftar ---
\setbeamercolor{itemize item}{fg=obeTeal}
\setbeamercolor{itemize subitem}{fg=obeTeal}
\setbeamercolor{enumerate item}{fg=obeBlue}
\setbeamertemplate{itemize item}{\raise0.2ex\hbox{\tiny$\blacksquare$}}
\setbeamertemplate{itemize subitem}{\raise0.2ex\hbox{\scriptsize$\bullet$}}

% --- kotak (pengganti tcolorbox pada buku) ---
\setbeamercolor{block title}{bg=obeBlueLight,fg=obeBlue}
\setbeamercolor{block body}{bg=obeBlueLight!45,fg=black!88}
\setbeamercolor{block title alerted}{bg=obeMaroonLight,fg=obeMaroon}
\setbeamercolor{block body alerted}{bg=obeMaroonLight!45,fg=black!88}
\setbeamercolor{block title example}{bg=obeTealLight,fg=obeTeal}
\setbeamercolor{block body example}{bg=obeTealLight!45,fg=black!88}

% ---------------------------------------------------------
% 5. Makro notasi (sama persis dengan buku)
%
% Catatan: \note TIDAK didefinisikan di sini karena beamer
% sudah memakai \note untuk catatan pembicara. Untuk catatan
% di slide dipakai \catatan.
% ---------------------------------------------------------
\newcommand{\code}[1]{\texttt{#1}}
\newcommand{\asm}[1]{\texttt{#1}}
\newcommand{\reg}[1]{\texttt{#1}}
\newcommand{\flag}[1]{\texttt{#1}}
\newcommand{\dir}[1]{\texttt{#1}}
\newcommand{\file}[1]{\texttt{#1}}
\newcommand{\eng}[1]{\textit{#1}}

\newcommand{\hex}[1]{\texttt{#1h}}
\newcommand{\bin}[1]{\texttt{#1b}}
\newcommand{\dec}[1]{\texttt{#1}}

\newcommand{\important}[1]{\textbf{Penting:} #1}
\newcommand{\catatan}[1]{\textbf{Catatan:} #1}
\newcommand{\term}[1]{\textbf{#1}}

\newcommand{\CPL}[1]{\textbf{CPL#1}}
\newcommand{\CPMK}[1]{\textbf{CPMK#1}}
\newcommand{\SubCPMK}[1]{\textbf{Sub-CPMK #1}}

\newcommand{\checkbox}{\ding{113}\hspace{0.35em}}
\newcommand{\checked}{\ding{51}\hspace{0.35em}}
\newcommand{\crossed}{\ding{55}\hspace{0.35em}}

\newcommand{\flowarrow}{\ensuremath{\rightarrow}}

% ---------------------------------------------------------
% 6. Makro tata letak slide
% ---------------------------------------------------------

% Keterangan sumber materi (ketertelusuran ke buku)
\newcommand{\Sumber}[1]{%
    \vfill
    {\scriptsize\color{obeGray}Sumber: #1\par}%
}

% Frame judul deck: \JudulDeck{<judul bab>}{<pertemuan>}
\newcommand{\JudulDeck}[2]{%
\begin{frame}[plain]
    \vspace{0.25cm}
    {\color{obeBlue}\sffamily\bfseries\small \NamaMK\ (\KodeMK)\par}
    \vspace{0.12cm}
    {\color{obeGray}\sffamily\scriptsize
     \Prodi\ \textbullet\ \Fakultas\ \textbullet\ \PerguruanTinggi\par}
    \vfill
    {\color{obeBlue}\sffamily\LARGE\bfseries #1\par}
    \vspace{0.4cm}
    {\color{obeTeal}\sffamily\large #2\par}
    \vspace{0.45cm}
    {\color{obeTeal}\rule{\textwidth}{1pt}}
    \vfill
    {\color{obeGray}\sffamily\small Dosen Pengampu\par}
    \vspace{0.08cm}
    {\sffamily\bfseries \DosenPengampu\par}
    \vspace{0.2cm}
    {\color{obeGray}\sffamily\small \SemesterTeks\par}
    \vspace{0.25cm}
\end{frame}
}

% Frame pemisah bagian
\newcommand{\PemisahBagian}[1]{%
\begin{frame}[plain]
    \vfill
    \begin{center}
        {\color{obeTeal}\sffamily\small\bfseries BAGIAN\par}
        \vspace{0.35cm}
        {\color{obeBlue}\sffamily\Large\bfseries #1\par}
        \vspace{0.35cm}
        {\color{obeTeal}\rule{0.55\textwidth}{0.9pt}}
    \end{center}
    \vfill
\end{frame}
}

% ---------------------------------------------------------
% 7. Gaya penampil kode
% ---------------------------------------------------------
% =========================================================
% ppt/src/listing-ppt.tex
% Gaya penampil kode untuk slide.
%
% Diturunkan dari buku_ajar/config/listings.tex, tetapi
% disesuaikan untuk slide: ukuran lebih kecil (karena frame
% 16:9 dan kode harus utuh terlihat) dan jarak lebih rapat.
% =========================================================

% ---------------------------------------------------------
% 1. Definisi bahasa MASM/TASM 8086 (salinan dari buku)
% ---------------------------------------------------------
\lstdefinelanguage{MASM}{
    morekeywords={
        MOV, XCHG, LEA, LDS, LES, PUSH, POP, PUSHA, POPA, PUSHF, POPF,
        LAHF, SAHF, XLAT, IN, OUT, CBW, CWD,
        ADD, ADC, SUB, SBB, INC, DEC, NEG, CMP, MUL, IMUL, DIV, IDIV,
        AAA, AAS, AAM, AAD, DAA, DAS,
        AND, OR, XOR, NOT, TEST, SHL, SHR, SAL, SAR, ROL, ROR, RCL, RCR,
        JMP, JE, JZ, JNE, JNZ, JG, JGE, JL, JLE, JA, JAE, JB, JBE,
        JC, JNC, JS, JNS, JO, JNO, JP, JPE, JNP, JPO, JCXZ,
        LOOP, LOOPE, LOOPZ, LOOPNE, LOOPNZ, CALL, RET, INT, IRET, INTO,
        HLT, NOP, CLI, STI, CLD, STD, CLC, STC, CMC, WAIT, ESC, LOCK,
        DB, DW, DD, DQ, DT, EQU, ORG, SEGMENT, ENDS, ASSUME, END,
        PROC, ENDP, MACRO, ENDM, LOCAL, PUBLIC, EXTRN, EXTERN,
        INCLUDE, INCLUDELIB, MODEL, STACK, CODE, DATA, CONST, ALIGN, EVEN,
        OFFSET, PTR, BYTE, WORD, DWORD, QWORD, TBYTE,
        NEAR, FAR, SHORT, LABEL, TYPE, SIZE, LENGTH, DUP, THIS,
        AX, BX, CX, DX, AL, AH, BL, BH, CL, CH, DL, DH,
        SP, BP, SI, DI, IP, CS, DS, SS, ES, FLAGS,
        CF, PF, AF, ZF, TF, IF, DF, OF
    },
    sensitive=false,
    morecomment=[l]{;},
    morestring=[b]',
    morestring=[b]"
}

% ---------------------------------------------------------
% 2. Gaya kode assembly untuk slide
% ---------------------------------------------------------
\lstdefinestyle{asmslide}{
    language=MASM,
    basicstyle=\ttfamily\scriptsize,
    keywordstyle=\color{obeBlue}\bfseries,
    commentstyle=\color{obeGreen}\itshape,
    stringstyle=\color{obeMaroon},
    numberstyle=\tiny\color{obeGray},
    numbers=left,
    numbersep=5pt,
    stepnumber=1,
    backgroundcolor=\color{codeBack},
    frame=single,
    rulecolor=\color{codeFrame},
    framesep=4pt,
    breaklines=true,
    breakatwhitespace=false,
    % CATATAN PENTING: opsi `postbreak` TIDAK boleh dipakai di sini.
    % Di beamer, `breaklines=true` bersama `postbreak=<kode>` membuat
    % listings gagal dirakit dengan galat
    %   "Missing number, treated as zero <to be read again> \kern".
    % Karena \lstset{style=asmslide} bersifat global, opsi postbreak
    % juga bocor ke gaya lain (mis. outslide/keluaran) dan merusaknya.
    % Di buku ajar (kelas `book`) justru sebaliknya: postbreak aman dan
    % dipakai. Jadi larangan ini khusus untuk slide.
    showstringspaces=false,
    tabsize=4,
    columns=fullflexible,
    keepspaces=true,
    upquote=true,
    xleftmargin=12pt,
    aboveskip=0.4em,
    belowskip=0.4em
}

% ---------------------------------------------------------
% 3. Gaya keluaran program / sesi debugger untuk slide
% ---------------------------------------------------------
\lstdefinestyle{outslide}{
    basicstyle=\ttfamily\scriptsize,
    keywordstyle=\color{black},
    commentstyle=\color{obeGray},
    stringstyle=\color{black},
    numbers=none,
    backgroundcolor=\color{obeGrayLight},
    frame=single,
    rulecolor=\color{codeFrame},
    framesep=4pt,
    breaklines=true,
    showstringspaces=false,
    columns=fullflexible,
    keepspaces=true,
    upquote=true,
    xleftmargin=4pt,
    aboveskip=0.4em,
    belowskip=0.4em
}

\lstset{style=asmslide}

% ---------------------------------------------------------
% 4. Environment ringkas untuk dipakai di dalam frame
%
% Frame yang memuat kedua environment ini WAJIB memakai
% opsi [fragile], karena isinya verbatim.
% ---------------------------------------------------------
\lstnewenvironment{kodeasm}[1][]{%
    \lstset{style=asmslide, #1}%
}{}

\lstnewenvironment{keluaran}[1][]{%
    \lstset{style=outslide, #1}%
}{}



\renewcommand{\PertemuanTeks}{Pertemuan 2}

\begin{document}

% =========================================================
% 1. JUDUL
% =========================================================
\JudulDeck{Arsitektur Mikroprosesor 8086}{Pertemuan 2}

% =========================================================
% 2. CAPAIAN DAN TUJUAN
% =========================================================
\begin{frame}{Capaian dan Tujuan Pembelajaran}
\begin{block}{Sub-CPMK 1.1}
\small Menjelaskan arsitektur dasar mikroprosesor 8086
(\eng{Execution Unit} dan \eng{Bus Interface Unit}), peran bus
alamat/data/kontrol, serta tahapan siklus instruksi
(\eng{fetch}--\eng{decode}--\eng{execute}).
\end{block}

\vspace{0.15em}
{\small\textbf{CPMK yang didukung:} \CPMK{01}\quad\textbullet\quad
\textbf{CPL yang didukung:} \CPL{02}\par}

\vspace{0.2em}
\begin{exampleblock}{Setelah pertemuan ini, mahasiswa mampu:}
\small
\begin{enumerate}
    \item menjelaskan kedudukan 8086 sebagai dasar keluarga x86
          beserta evolusinya;
    \item membedakan arsitektur internal dan arsitektur eksternal
          mikroprosesor;
    \item menguraikan pembagian tugas BIU dan EU serta peran
          \eng{instruction queue};
    \item menyebutkan fungsi bus alamat, data, dan kontrol beserta
          sinyal utamanya;
    \item menelusuri satu instruksi melalui
          \eng{fetch}--\eng{decode}--\eng{execute}--\eng{store}.
\end{enumerate}
\end{exampleblock}
\end{frame}

% =========================================================
% 3. PETA MATERI
% =========================================================
\begin{frame}{Peta Materi}
\small
\begin{enumerate}
    \item \textbf{Evolusi Mikroprosesor dan Kedudukan 8086}\\
          {\footnotesize dari 4004 sampai 8086; mengapa 8086 masih
          dipelajari; arsitektur internal dan eksternal}

    \vspace{0.25em}
    \item \textbf{Bus Interface Unit dan Execution Unit}\\
          {\footnotesize diagram blok 8086; pembagian tugas unit;
          \eng{instruction queue} dan \eng{pipelining}; alamat fisik}

    \vspace{0.25em}
    \item \textbf{Bus Eksternal dan Siklus Instruksi}\\
          {\footnotesize bus alamat, data, kontrol;
          \eng{multiplexing} AD0--AD15; siklus
          \eng{fetch}--\eng{store}}
\end{enumerate}

\vspace{0.4em}
\begin{block}{Praktik pada pertemuan ini}
\small Menggambar diagram blok BIU/EU dari ingatan, lalu menelusuri
siklus instruksi dengan \eng{Single Step} pada Emu8086.
\end{block}
\end{frame}

% =========================================================
% BAGIAN 1
% =========================================================
\PemisahBagian{Evolusi Mikroprosesor dan Kedudukan 8086}

% ---------------------------------------------------------
\begin{frame}{Dari 4004 ke 8086}
\begin{center}
\footnotesize
\begin{tabular}{@{}p{0.95cm}p{1.75cm}p{1.25cm}p{1.35cm}p{6.9cm}@{}}
    \toprule
    \textbf{Tahun} & \textbf{Prosesor} & \textbf{Bus data} &
    \textbf{Bus alamat} & \textbf{Catatan} \\
    \midrule
    1971 & Intel 4004 & 4 bit & 12 bit &
    Mikroprosesor komersial pertama, untuk kalkulator \\
    \addlinespace
    1974 & Intel 8080 & 8 bit & 16 bit &
    Ruang alamat 64 KB; \eng{microcomputer} awal \\
    \addlinespace
    1978 & Intel 8086 & 16 bit & 20 bit &
    Ruang alamat 1 MB; dasar keluarga x86 \\
    \addlinespace
    1979 & Intel 8088 & 8 bit (luar) & 20 bit &
    Inti 16-bit, bus luar 8-bit; dipakai IBM PC \\
    \addlinespace
    1982 & Intel 80286 & 16 bit & 24 bit &
    Menambah \eng{protected mode}, 16 MB alamat \\
    \bottomrule
\end{tabular}
\end{center}

\vspace{0.35em}
{\footnotesize\color{obeGray}Setiap generasi menambah lebar kata,
lebar bus alamat, dan jumlah transistor. Varian 8088 berinti 16 bit
hanya menyempitkan bus luarnya menjadi 8 bit.\par}

\Sumber{Bab IV, Bagian 1.1 --- Tabel 4.1}
\end{frame}

% ---------------------------------------------------------
\begin{frame}{Mengapa 8086 Masih Dipelajari}
\small
\begin{enumerate}
    \item \textbf{Sederhana namun lengkap.}
          Sudah memiliki seluruh konsep penting prosesor modern:
          register, bus, \eng{flag}, segmentasi, dan
          \eng{pipeline} sederhana, tanpa mekanisme rumit seperti
          \eng{cache} berlapis.

    \vspace{0.35em}
    \item \textbf{Mudah diamati.}
          Emu8086 memperlihatkan isi register, memori, dan
          \eng{flag} pada setiap langkah eksekusi, sehingga konsep
          abstrak dapat diverifikasi sendiri.

    \vspace{0.35em}
    \item \textbf{Akar keluarga x86.}
          Prosesor x86 modern tetap kompatibel ke belakang; program
          8086 mode real masih dapat dijalankan melalui mode
          kompatibilitas.
\end{enumerate}

\vspace{0.45em}
\begin{alertblock}{Intinya}
\small Usianya lebih dari empat dekade, tetapi 8086 adalah model
paling jernih untuk memahami cara kerja sebuah prosesor.
\end{alertblock}
\end{frame}

% ---------------------------------------------------------
\begin{frame}{Arsitektur Internal dan Eksternal}
\begin{exampleblock}{Arsitektur internal (\eng{internal architecture})}
\small Susunan unit fungsional \textbf{di dalam chip}: register,
ALU, unit pengendali, dan saluran data internal. Menjelaskan
\emph{bagaimana} prosesor mengolah sebuah instruksi.
\end{exampleblock}

\vspace{0.35em}
\begin{block}{Arsitektur eksternal (\eng{external architecture})}
\small Tampak \textbf{luar chip}: jumlah dan fungsi pin, kelompok
bus, serta sinyal kendali ke memori dan perangkat I/O.
Menjelaskan \emph{bagaimana} prosesor berbicara dengan dunia luar.
\end{block}

\vspace{0.45em}
{\footnotesize\color{obeGray}Programmer assembly lebih sering
berhadapan dengan arsitektur internal (menuliskan nama register);
perancang sistem digital lebih banyak bekerja pada arsitektur
eksternal (pin \reg{A0}--\reg{A19}, \reg{AD0}--\reg{AD15}, sinyal
kendali).\par}
\end{frame}

% ---------------------------------------------------------
\begin{frame}{Cara Pandang Sederhana: Kotak Hitam}
\vspace{0.2em}
\begin{center}
\begin{tikzpicture}[
    kotak/.style={
        draw=obeBlue, line width=1pt, fill=obeBlueLight,
        rounded corners=3pt, minimum width=3.8cm, minimum height=1.5cm,
        align=center, font=\small\bfseries
    },
    bus/.style={
        draw=obeTeal, line width=0.8pt, fill=obeTealLight,
        rounded corners=2pt, minimum width=4.4cm, minimum height=0.62cm,
        align=center, font=\scriptsize
    },
    panah/.style={line width=0.8pt, obeTeal, -{Stealth[length=2mm]}},
    panah2/.style={
        line width=0.8pt, obeTeal,
        {Stealth[length=2mm]}-{Stealth[length=2mm]}
    }
]
    \node[kotak] (cpu) at (0,0) {Mikroprosesor\\8086};
    \node[bus, right=1.1cm of cpu, yshift=1.35cm] (a)
        {Bus alamat 20-bit};
    \node[bus, right=1.1cm of cpu, yshift=0.45cm] (d)
        {Bus data 16-bit};
    \node[bus, right=1.1cm of cpu, yshift=-0.45cm] (c)
        {Bus kontrol};
    \node[bus, right=1.1cm of cpu, yshift=-1.35cm] (p)
        {Catuan \& pewaktu};
    \draw[panah] (cpu.east) -- (a.west);
    \draw[panah2] (cpu.east) -- (d.west);
    \draw[panah2] (cpu.east) -- (c.west);
    \draw[panah] (cpu.east) -- (p.west);
\end{tikzpicture}
\end{center}

\vspace{0.5em}
\small Ke luar, prosesor hanya berhubungan melalui empat kelompok
saluran tersebut. Setiap instruksi yang dijalankan pasti melibatkan
pertukaran informasi lewat salah satunya.
\end{frame}

% ---------------------------------------------------------
\begin{frame}{Ringkasan Bagian 1}
\small
\begin{itemize}
    \item Mikroprosesor berkembang dari Intel 4004 (4 bit) sampai
          8086 (16 bit; bus alamat 20 bit) yang menjadi dasar
          keluarga x86.

    \vspace{0.3em}
    \item Varian 8088 berinti 16 bit dengan bus luar 8 bit dipakai
          pada IBM PC.

    \vspace{0.3em}
    \item 8086 dipelajari karena sederhana namun lengkap, mudah
          diamati pada emulator, dan menjadi akar keluarga x86.

    \vspace{0.3em}
    \item Arsitektur dipandang dari dua sisi: \term{internal}
          (register, ALU, unit di dalam chip) dan \term{eksternal}
          (pin, bus, sinyal kendali).
\end{itemize}
\end{frame}

% =========================================================
% BAGIAN 2
% =========================================================
\PemisahBagian{Bus Interface Unit dan Execution Unit}

% ---------------------------------------------------------
\begin{frame}{Diagram Blok Internal 8086}
\begin{center}
\begin{tikzpicture}[
    unit/.style={
        draw=obeBlue, line width=1pt, rounded corners=3pt,
        fill=obeBlueLight, minimum width=4.75cm, minimum height=2.7cm
    },
    unitb/.style={
        draw=obeTeal, line width=1pt, rounded corners=3pt,
        fill=obeTealLight, minimum width=4.75cm, minimum height=2.7cm
    },
    isi/.style={
        draw=obeBlue, fill=white, minimum width=4.0cm,
        minimum height=0.5cm, align=center, font=\scriptsize
    },
    isib/.style={
        draw=obeTeal, fill=white, minimum width=4.0cm,
        minimum height=0.5cm, align=center, font=\scriptsize
    },
    judul/.style={font=\small\bfseries},
    panah/.style={
        {Stealth[length=2.2mm]}-{Stealth[length=2.2mm]},
        line width=1.3pt, obeBlue
    }
]
    % ---------------- EU ----------------
    \node[unit] (eu) {};
    \node[judul, text=obeBlue, anchor=north]
        at ([yshift=-4pt]eu.north) {Execution Unit (EU)};
    \node[isi] (e1) at ([yshift=-0.85cm]eu.north) {ALU};
    \node[isi, below=0.14cm of e1] (e2)
        {Register umum, penunjuk, indeks};
    \node[isi, below=0.14cm of e2] (e3) {Register \eng{flag}};

    % ---------------- BIU ----------------
    \node[unitb, right=1.55cm of eu] (biu) {};
    \node[judul, text=obeTeal, anchor=north]
        at ([yshift=-4pt]biu.north) {Bus Interface Unit (BIU)};
    \node[isib] (b1) at ([yshift=-0.85cm]biu.north)
        {Register segmen \& \reg{IP}};
    \node[isib, below=0.14cm of b1] (b2)
        {\eng{Instruction queue} 6 byte};
    \node[isib, below=0.14cm of b2] (b3)
        {Penghitung alamat 20-bit};

    \draw[panah] (eu.east) -- (biu.west)
        node[midway, above, font=\tiny, text=obeBlue] {16-bit};

    % ---------------- bus luar ----------------
    \node[draw=obeAmber, fill=obeAmberLight, rounded corners=2pt,
          minimum width=4.3cm, minimum height=0.62cm, align=center,
          font=\scriptsize, below=0.95cm of biu] (ext)
          {Bus eksternal: alamat 20-bit / data 16-bit};
    \draw[panah, obeAmber] (biu.south) -- (ext.north);
    \node[draw=obeGray, fill=obeGrayLight, rounded corners=2pt,
          minimum width=4.3cm, minimum height=0.62cm, align=center,
          font=\scriptsize, below=0.4cm of ext] (mem)
          {Memori dan perangkat I/O};
    \draw[panah, obeGray] (ext.south) -- (mem.north);
\end{tikzpicture}
\end{center}

\vspace{0.15em}
{\footnotesize\color{obeGray}EU di kiri, BIU di kanan, dihubungkan
bus internal 16 bit. Hanya BIU yang menyentuh pin luar chip.\par}
\end{frame}

% ---------------------------------------------------------
\begin{frame}{Tugas Masing-masing Unit}
\begin{center}
\footnotesize
\begin{tabular}{@{}p{2.4cm}p{5.3cm}p{5.3cm}@{}}
    \toprule
    \textbf{Aspek} & \textbf{Bus Interface Unit (BIU)} &
    \textbf{Execution Unit (EU)} \\
    \midrule
    Tugas utama & Mengambil instruksi dan data; menangani seluruh
    lalu lintas bus & Menerjemahkan dan menjalankan instruksi \\
    \addlinespace
    Register & \reg{CS}, \reg{DS}, \reg{SS}, \reg{ES}, \reg{IP},
    \eng{instruction queue} & \reg{AX}--\reg{DX}, \reg{SP},
    \reg{BP}, \reg{SI}, \reg{DI}, \eng{flag} \\
    \addlinespace
    Akses bus luar & Wewenang tunggal &
    Tidak punya; harus meminta melalui BIU \\
    \addlinespace
    Perhitungan alamat & Alamat fisik 20-bit dari pasangan
    segmen:offset & Hanya \eng{offset} (\eng{effective address}) \\
    \bottomrule
\end{tabular}
\end{center}

\vspace{0.4em}
\begin{alertblock}{Yang perlu diingat}
\small Ketika sebuah instruksi memerlukan operand dari memori, EU
mengirim permintaan alamat ke BIU, lalu BIU menjalankan siklus bus
dan menyerahkan datanya kembali.
\end{alertblock}
\end{frame}

% ---------------------------------------------------------
\begin{frame}{Instruction Queue dan Pipelining}
\small Selama EU masih sibuk, BIU memakai waktu luang bus untuk
mengisi \eng{instruction queue} enam byte.

\vspace{0.35em}
\begin{center}
\begin{tikzpicture}[x=1.4cm, y=0.55cm, font=\scriptsize]
    % --- baris 1: tanpa antrean ---
    \node[anchor=east, font=\scriptsize\bfseries] at (-0.1,2.6)
        {Tanpa antrean};
    \foreach \i/\lbl/\c in {%
        0/F1/obeBlueLight, 1/E1/obeGreenLight,
        2/F2/obeBlueLight, 3/E2/obeGreenLight,
        4/F3/obeBlueLight, 5/E3/obeGreenLight} {
        \draw[draw=obeGray, fill=\c, line width=0.4pt]
            (\i,2.2) rectangle (\i+1,3.0);
        \node at (\i+0.5,2.6) {\lbl};
    }
    % --- baris 2: dengan antrean ---
    \node[anchor=east, font=\scriptsize\bfseries] at (-0.1,1.1)
        {Dengan antrean};
    \draw[draw=obeGray, fill=obeBlueLight, line width=0.4pt]
        (0,0.4) rectangle (1,1.8);
    \node at (0.5,1.1) {F1};
    \draw[draw=obeGray, fill=obeGreenLight, line width=0.4pt]
        (1,1.1) rectangle (2,1.8);
    \node[font=\tiny] at (1.5,1.45) {E1};
    \draw[draw=obeGray, fill=obeBlueLight, line width=0.4pt]
        (1,0.4) rectangle (2,1.1);
    \node[font=\tiny] at (1.5,0.75) {F2};
    \draw[draw=obeGray, fill=obeGreenLight, line width=0.4pt]
        (2,1.1) rectangle (3,1.8);
    \node[font=\tiny] at (2.5,1.45) {E2};
    \draw[draw=obeGray, fill=obeBlueLight, line width=0.4pt]
        (2,0.4) rectangle (3,1.1);
    \node[font=\tiny] at (2.5,0.75) {F3};
    \draw[draw=obeGray, fill=obeGreenLight, line width=0.4pt]
        (3,0.4) rectangle (4,1.8);
    \node at (3.5,1.1) {E3};
    % --- sumbu siklus ---
    \foreach \i in {1,...,6} {
        \node[font=\tiny, text=obeGray] at (\i-0.5,0.05) {\i};
    }
    \node[font=\tiny, text=obeGray, anchor=west] at (6.05,0.05)
        {siklus};
\end{tikzpicture}
\end{center}

\vspace{0.2em}
{\footnotesize F = \eng{fetch} (BIU) \quad\textbullet\quad
E = \eng{execute} (EU) \quad\textbullet\quad Tanpa antrean: 6
siklus; dengan antrean: 4 siklus.\par}
\end{frame}

% ---------------------------------------------------------
\begin{frame}{Peringatan: Queue Flush}
\begin{alertblock}{\eng{Pipeline} pada 8086 hanya satu tingkat}
\small Antreannya hanya enam byte. Ketika EU menjalankan instruksi
lompatan (\asm{JMP}, \asm{CALL}), isi antrean sudah tidak relevan
sehingga harus dikosongkan. Peristiwa ini disebut
\term{queue flush}.
\end{alertblock}

\vspace{0.5em}
\small
\begin{itemize}
    \item \eng{Queue flush} menghapus keuntungan \eng{pipeline}
          untuk instruksi tersebut.

    \vspace{0.25em}
    \item Akibatnya, program yang penuh lompatan berjalan lebih
          lambat daripada program lurus dengan panjang yang sama.
\end{itemize}

\Sumber{Bab IV, Bagian 2.3}
\end{frame}

% ---------------------------------------------------------
\begin{frame}{Peran Segmentasi pada BIU}
\small BIU mengubah pasangan \term{segmen:offset} menjadi alamat
fisik 20-bit. Register segmen hanya 16 bit, sedangkan bus alamat
selebar 20 bit.

\vspace{0.5em}
\begin{center}
{\color{obeBlue}\large Alamat fisik $=$ (register segmen $\times$
16) $+$ \eng{offset}}
\end{center}

\vspace{0.5em}
\begin{exampleblock}{Contoh}
\small \reg{CS} $=$ \hex{1234}, \reg{IP} $=$ \hex{0010}\\[0.25em]
\hex{12340} $+$ \hex{0010} $=$ \textbf{\hex{12350}}\\[0.25em]
{\footnotesize Mengalikan dengan 16 sama dengan menggeser nilai
segmen satu digit heksadesimal ke kiri.}
\end{exampleblock}

\vspace{0.35em}
{\footnotesize\color{obeGray}Rumus inilah alasan 8086 mampu
menjangkau 1 MB memori meskipun register-nya hanya 16 bit.\par}
\end{frame}

% ---------------------------------------------------------
\begin{frame}{Sekilas Register Internal}
\begin{center}
\footnotesize
\begin{tabular}{@{}p{2.7cm}p{3.1cm}p{6.7cm}@{}}
    \toprule
    \textbf{Kelompok} & \textbf{Anggota} & \textbf{Kegunaan utama} \\
    \midrule
    Register umum & \reg{AX}, \reg{BX}, \reg{CX}, \reg{DX} &
    Menyimpan data dan hasil operasi aritmetika \\
    \addlinespace
    Register penunjuk & \reg{SP}, \reg{BP} &
    Puncak \eng{stack} dan kerangka data prosedur \\
    \addlinespace
    Register indeks & \reg{SI}, \reg{DI} &
    Sumber dan tujuan pada operasi string \\
    \addlinespace
    Register segmen & \reg{CS}, \reg{DS}, \reg{SS}, \reg{ES} &
    Nomor segmen kode, data, \eng{stack}, dan tambahan \\
    \addlinespace
    Penunjuk instruksi & \reg{IP} &
    \eng{Offset} instruksi berikutnya \\
    \addlinespace
    Register \eng{flag} & \flag{CF}, \flag{ZF}, \flag{SF}, \dots &
    Kondisi hasil operasi terakhir \\
    \bottomrule
\end{tabular}
\end{center}

\vspace{0.35em}
{\footnotesize\color{obeGray}Register segmen, \reg{IP}, dan
\eng{instruction queue} berada di BIU; sisanya berada di EU.
Rincian setiap register dibahas pada Bab V.\par}
\end{frame}

% ---------------------------------------------------------
\begin{frame}{Ringkasan Bagian 2}
\small
\begin{itemize}
    \item 8086 terbagi menjadi BIU dan EU yang bekerja secara
          paralel.

    \vspace{0.3em}
    \item BIU: register segmen, \reg{IP}, penghitung alamat 20-bit,
          dan \eng{instruction queue} enam byte; pemegang tunggal
          bus luar.

    \vspace{0.3em}
    \item EU: ALU, register umum/penunjuk/indeks, dan register
          \eng{flag}; sama sekali tidak terhubung ke pin luar.

    \vspace{0.3em}
    \item Antrean memungkinkan \eng{pipelining} (6 siklus menjadi 4
          siklus); alamat fisik $=$ (segmen $\times$ 16) $+$
          \eng{offset}.
\end{itemize}
\end{frame}

% =========================================================
% BAGIAN 3
% =========================================================
\PemisahBagian{Bus Eksternal dan Siklus Instruksi}

% ---------------------------------------------------------
\begin{frame}{Tiga Kelompok Bus}
\small Bus adalah sekumpulan penghantar yang dipakai bersama oleh
prosesor, memori, dan perangkat I/O.

\vspace{0.35em}
\begin{itemize}
    \item \textbf{Bus alamat} (\eng{address bus}) --- 20 bit, yaitu
          \reg{A0}--\reg{A19}; membawa alamat lokasi yang akan
          diakses. Menjangkau $2^{20}$ byte $=$ 1 MB dan bersifat
          \textbf{satu arah} (prosesor $\rightarrow$ memori/I/O).

    \vspace{0.3em}
    \item \textbf{Bus data} (\eng{data bus}) --- 16 bit, yaitu
          \reg{AD0}--\reg{AD15}; membawa data yang sedang dibaca
          atau ditulis, dan bersifat \textbf{dua arah}.

    \vspace{0.3em}
    \item \textbf{Bus kontrol} (\eng{control bus}) --- mengatur
          \emph{jenis} dan \emph{arah} pertukaran data: baca atau
          tulis, memori atau perangkat I/O, dan kapan perangkat
          siap melayani.
\end{itemize}
\end{frame}

% ---------------------------------------------------------
\begin{frame}{\eng{Multiplexing} \reg{AD0}--\reg{AD15}}
\small Kemasan 8086 hanya \textbf{40 pin}, sehingga sebagian pin
dipakai bergantian: \reg{AD0}--\reg{AD15} membawa alamat pada awal
siklus bus, lalu membawa data.

\vspace{0.6em}
\begin{center}
\begin{tikzpicture}[font=\scriptsize]
    \draw[draw=obeBlue, fill=obeBlueLight, line width=0.9pt]
        (0,0) rectangle (5,0.95);
    \node at (2.5,0.475) {Alamat \reg{A0}--\reg{A15}};
    \draw[draw=obeTeal, fill=obeTealLight, line width=0.9pt]
        (5,0) rectangle (10,0.95);
    \node at (7.5,0.475) {Data \reg{D0}--\reg{D15}};

    \draw[{Stealth[length=2mm]}-, line width=0.9pt, obeAmber]
        (5,1.5) -- (5,1.0);
    \node[anchor=south, font=\scriptsize, text=obeAmber]
        at (5,1.5) {\reg{ALE} $=1$: tangkap alamat};
    \node[anchor=north, font=\scriptsize, text=obeGray]
        at (0,0) {awal siklus};
    \node[anchor=north, font=\scriptsize, text=obeGray]
        at (10,0) {akhir siklus};
\end{tikzpicture}
\end{center}

\vspace{0.5em}
{\footnotesize\color{obeGray}Karena alamat hanya muncul sebentar,
rangkaian luar bernama \eng{latch} harus menangkapnya. Sinyal
\reg{ALE} (\eng{Address Latch Enable}) memberi tahu kapan saat yang
tepat.\par}
\end{frame}

% ---------------------------------------------------------
\begin{frame}{Sinyal Kendali Utama 8086}
\begin{center}
\footnotesize
\begin{tabular}{@{}p{1.4cm}p{1.6cm}p{9.8cm}@{}}
    \toprule
    \textbf{Sinyal} & \textbf{Arah} & \textbf{Fungsi} \\
    \midrule
    \reg{ALE} & Keluaran &
    Menandai bahwa \reg{AD0}--\reg{AD15} sedang membawa alamat \\
    \addlinespace
    \code{M/IO} & Keluaran &
    Sasaran akses: tinggi untuk memori, rendah untuk perangkat I/O \\
    \addlinespace
    \reg{RD} & Keluaran &
    Aktif rendah; prosesor sedang melakukan operasi baca \\
    \addlinespace
    \reg{WR} & Keluaran &
    Aktif rendah; prosesor sedang melakukan operasi tulis \\
    \addlinespace
    \reg{READY} & Masukan &
    Perangkat lambat sudah selesai; bila rendah, prosesor menambah
    \eng{wait state} \\
    \addlinespace
    \reg{RESET} & Masukan &
    Kembali ke keadaan awal; \reg{CS} diisi \hex{FFFF} dan
    \reg{IP} diisi \hex{0000} \\
    \bottomrule
\end{tabular}
\end{center}

\vspace{0.3em}
{\footnotesize\color{obeGray}Sinyal \code{MN/MX} memilih mode
minimum atau maksimum. Seluruh contoh pada buku ini dapat dipandang
berjalan pada mode minimum.\par}
\end{frame}

% ---------------------------------------------------------
\begin{frame}{Siklus Instruksi 8086}
\vspace{0.2em}
\begin{center}
\begin{tikzpicture}[
    tahap/.style={
        draw=obeBlue, line width=0.9pt, fill=obeBlueLight,
        rounded corners=3pt, minimum width=2.25cm,
        minimum height=1.1cm, align=center, font=\small\bfseries
    },
    panah/.style={-{Stealth[length=2.2mm]}, line width=0.9pt, obeTeal},
    unit/.style={font=\scriptsize, text=obeGray}
]
    \node[tahap] (f) {Fetch\\{\normalfont\scriptsize ambil opcode}};
    \node[tahap, right=0.4cm of f, fill=obeTealLight, draw=obeTeal] (d)
        {Decode\\{\normalfont\scriptsize baca opcode}};
    \node[tahap, right=0.4cm of d, fill=obeAmberLight, draw=obeAmber] (e)
        {Execute\\{\normalfont\scriptsize jalankan}};
    \node[tahap, right=0.4cm of e, fill=obeMaroonLight, draw=obeMaroon]
        (s) {Store\\{\normalfont\scriptsize simpan hasil}};

    \draw[panah] (f) -- (d);
    \draw[panah] (d) -- (e);
    \draw[panah] (e) -- (s);

    \node[unit, below=0.08cm of f] {BIU};
    \node[unit, below=0.08cm of d] {EU};
    \node[unit, below=0.08cm of e] {EU};
    \node[unit, below=0.08cm of s] {BIU};

    \draw[panah, dashed]
        (s.east) -- ([xshift=0.45cm]s.east)
        -- ([xshift=0.45cm, yshift=-1.5cm]s.east)
        -- ([xshift=-0.45cm, yshift=-1.5cm]f.west)
        -- ([xshift=-0.45cm]f.west) -- (f.west);
    \node[font=\scriptsize, text=obeGray]
        at ($([xshift=0.45cm]s.east)!0.5!([xshift=-0.45cm]f.west)
             + (0,-1.72cm)$)
        {instruksi berikutnya};
\end{tikzpicture}
\end{center}

\vspace{0.35em}
\small \term{Fetch}: BIU menghitung alamat dari \reg{CS}:\reg{IP}
dan mengisi antrean \quad\textbullet\quad \term{Decode}: EU
mengenali opcode \quad\textbullet\quad \term{Store}: hanya
memerlukan siklus bus bila tujuannya memori.
\end{frame}

% ---------------------------------------------------------
\begin{frame}{Menelusuri Satu Instruksi}
\small Instruksi \code{MOV AX, 1234h} berukuran tiga byte dengan
opcode \hex{B8} diikuti \hex{34} dan \hex{12} secara
\eng{little endian}. Andaikan \reg{CS} $=$ \hex{1000}, \reg{IP} $=$
\hex{0000}.

\vspace{0.35em}
\begin{center}
\footnotesize
\begin{tabular}{@{}p{0.7cm}p{1.4cm}p{1.5cm}p{1.1cm}p{7.3cm}@{}}
    \toprule
    \textbf{No.} & \textbf{Tahap} & \textbf{Alamat} &
    \textbf{\reg{IP}} & \textbf{Kejadian} \\
    \midrule
    1 & Fetch & \hex{10000} & \hex{0003} &
    BIU membaca \hex{B8}, \hex{34}, \hex{12}; antrean: \hex{B8 34 12} \\
    \addlinespace
    2 & Decode & --- & \hex{0003} &
    EU mengenali opcode: \code{MOV reg16, imm16} \\
    \addlinespace
    3 & Execute & --- & \hex{0003} &
    EU menyusun \hex{1234}, lalu menulis ke \reg{AX} \\
    \addlinespace
    4 & Store & --- & \hex{0003} &
    Tujuan register: tidak ada siklus bus \\
    \bottomrule
\end{tabular}
\end{center}

\vspace{0.35em}
{\footnotesize\color{obeGray}\reg{IP} bertambah tiga --- panjang
instruksi menentukan kenaikannya. Bandingkan \code{MOV hasil1, AX}
yang tujuannya di memori: BIU harus menjalankan siklus bus
tulis.\par}

\Sumber{Bab IV, Bagian 3.3 --- Tabel 4.6}
\end{frame}

% ---------------------------------------------------------
\begin{frame}{Ringkasan Bagian 3}
\small
\begin{itemize}
    \item Tiga kelompok bus: alamat 20-bit (satu arah), data 16-bit
          (dua arah), dan kontrol.

    \vspace{0.3em}
    \item Pin \reg{AD0}--\reg{AD15} dimultipleks; alamat ditangkap
          \eng{latch} dengan bantuan sinyal \reg{ALE}.

    \vspace{0.3em}
    \item Siklus instruksi: \eng{fetch} (BIU) $\rightarrow$
          \eng{decode} (EU) $\rightarrow$ \eng{execute} (EU)
          $\rightarrow$ \eng{store} (BIU).

    \vspace{0.3em}
    \item Tahap \eng{store} hanya memerlukan siklus bus bila
          tujuannya berada di memori.
\end{itemize}
\end{frame}

% =========================================================
% CONTOH PROGRAM
% =========================================================
\begin{frame}[fragile]{Contoh 1: Menelusuri Siklus Instruksi}
\small Program sesingkat mungkin agar setiap langkah dapat diamati
tanpa kehilangan jejak.

\vspace{0.3em}
\begin{columns}[t]
\begin{column}{0.5\textwidth}
\begin{kodeasm}
.MODEL SMALL
.STACK 100h
.DATA
    nilai_awal DW 1234h
.CODE
START:
    MOV AX, @DATA
    MOV DS, AX
    MOV AX, 1234h    ; B8 34 12
    MOV AH, 4Ch
    MOV AL, 00h
    INT 21h
END START
\end{kodeasm}
\end{column}
\begin{column}{0.46\textwidth}
{\color{obeTeal}\bfseries Panjang instruksi}\\[0.3em]
\begin{keluaran}
1000:0000  MOV AX, @DATA  3
1000:0003  MOV DS, AX     2
1000:0005  MOV AX, 1234h  3
1000:0008  MOV AH, 4Ch    2
1000:000A  MOV AL, 00h    2
1000:000C  INT 21h        2
\end{keluaran}

\vspace{0.35em}
{\footnotesize
\reg{IP} bertambah menurut panjang instruksi: 3 byte $\rightarrow$
$+3$, 2 byte $\rightarrow$ $+2$.\\[0.25em]
Dua instruksi 2 byte dapat berada dalam antrean enam byte
sekaligus.\\[0.25em]
Instruksi yang membaca variabel memicu siklus bus tambahan.\par}

\vspace{0.3em}
{\footnotesize\color{obeGray}\file{code/asm/siklus-dasar.asm}\par}
\end{column}
\end{columns}
\end{frame}

% ---------------------------------------------------------
\begin{frame}[fragile]{Contoh 2: Dua Cara Menjumlahkan}
\begin{columns}[t]
\begin{column}{0.55\textwidth}
\begin{kodeasm}
.CODE
START:
    MOV AX, @DATA
    MOV DS, AX
    ; cara 1: operand di register
    MOV AX, 1234h
    MOV BX, 5678h
    ADD AX, BX
    MOV hasil1, AX
    ; cara 2: operand di memori
    MOV AX, bil1
    ADD AX, bil2
    MOV hasil2, AX
    MOV AH, 4Ch
    INT 21h
\end{kodeasm}
\end{column}
\begin{column}{0.42\textwidth}
{\footnotesize\color{obeGray}\dir{.DATA}: \code{bil1} \hex{1234},
\code{bil2} \hex{5678}, \code{hasil1} \hex{0}, \code{hasil2}
\hex{0}.\par}

\vspace{0.4em}
{\color{obeTeal}\bfseries Keluaran}\\[0.3em]
\begin{keluaran}
hasil1 = 68ACh
hasil2 = 68ACh
\end{keluaran}

\vspace{0.4em}
{\footnotesize
Cara 1 (register): \textbf{11 siklus}.\\[0.25em]
Cara 2 (memori): \textbf{29 siklus}.\\[0.25em]
Penjumlahannya sama; yang bertambah adalah dua siklus bus pada
tahap \eng{store}.\par}

\vspace{0.3em}
{\footnotesize\color{obeGray}\file{code/asm/jumlah-dua-cara.asm}\par}
\end{column}
\end{columns}
\end{frame}

% =========================================================
% PRAKTIKUM
% =========================================================
\begin{frame}{Praktikum: Diagram Blok dan Penelusuran}
\begin{block}{Judul praktikum}
\small Menggambar diagram blok 8086 dan menelusuri siklus instruksi
dengan \eng{Single Step} \quad\textbullet\quad \textbf{Sub-CPMK:}
1.1 \quad\textbullet\quad \textbf{Waktu:} 100 menit
\quad\textbullet\quad \textbf{Bentuk:} individu
\end{block}

\vspace{0.3em}
\begin{columns}[t]
\begin{column}{0.48\textwidth}
{\color{obeBlue}\bfseries Alat dan bahan}
\begin{itemize}\small
    \item Komputer dengan Emu8086 terpasang
    \item \file{siklus-dasar.asm} dan
          \file{jumlah-dua-cara.asm}
    \item Kertas kosong dan alat tulis
    \item Lembar jurnal praktikum
\end{itemize}
\end{column}
\begin{column}{0.48\textwidth}
{\color{obeBlue}\bfseries Langkah kerja}
\begin{enumerate}\small
    \item Gambar diagram blok 8086 dari ingatan, lalu koreksi
    \item Catat \reg{CS}, \reg{IP}, \reg{AX}, \reg{DS} mula-mula
    \item Tekan \eng{Single Step}, catat perubahan tiap langkah
    \item Bandingkan panjang instruksi dengan hasil
          \eng{disassembly}
\end{enumerate}
\end{column}
\end{columns}

\vspace{0.35em}
{\footnotesize\textbf{Dikumpulkan:} gambar diagram blok beserta
koreksinya, tabel pengamatan \reg{CS}/\reg{IP}/antrean, dan
perbandingan jumlah siklus kedua cara penjumlahan.\par}
\end{frame}

% =========================================================
% LATIHAN
% =========================================================
\begin{frame}{Latihan}
\small
\begin{enumerate}
    \item Jelaskan perbedaan arsitektur internal dan arsitektur
          eksternal sebuah mikroprosesor. Sebutkan dua contoh
          masing-masing pada 8086.

    \vspace{0.3em}
    \item Mengapa 8086 dibagi menjadi BIU dan EU? Jelaskan
          keuntungannya dibandingkan prosesor dengan satu unit
          pengolah.

    \vspace{0.3em}
    \item Diketahui \reg{CS} $=$ \hex{2A10} dan \reg{IP} $=$
          \hex{00F3}. Hitung alamat fisik instruksi berikutnya dan
          tunjukkan langkah perhitungannya.

    \vspace{0.3em}
    \item Jelaskan mengapa instruksi \asm{JMP} menghapus keuntungan
          \eng{instruction queue}. Kaitkan dengan istilah
          \eng{queue flush}.

    \vspace{0.3em}
    \item Di Emu8086, \reg{IP} bertambah tiga setelah
          \code{MOV AX, 1234h} tetapi hanya bertambah satu setelah
          \code{INC AX}. Benarkah programnya salah? Jelaskan.
\end{enumerate}
\end{frame}

% =========================================================
% RANGKUMAN
% =========================================================
\begin{frame}{Rangkuman}
\small
\begin{itemize}
    \item Mikroprosesor berkembang dari Intel 4004 (4 bit) hingga
          8086 (16 bit) yang menjadi dasar keluarga x86; varian 8088
          dipakai pada IBM PC.

    \vspace{0.22em}
    \item Arsitektur dipandang dari dua sisi: \term{internal}
          (register, ALU, unit dalam chip) dan \term{eksternal}
          (pin, bus, sinyal kendali).

    \vspace{0.22em}
    \item 8086 terbagi menjadi BIU dan EU. BIU memegang seluruh lalu
          lintas bus luar; EU hanya menjalankan instruksi.

    \vspace{0.22em}
    \item Alamat fisik $=$ (segmen $\times$ 16) $+$ \eng{offset},
          sehingga 8086 menjangkau 1 MB meski register-nya 16 bit.

    \vspace{0.22em}
    \item \eng{Instruction queue} enam byte memungkinkan
          \eng{pipelining}; lompatan mengosongkan antrean
          (\eng{queue flush}).

    \vspace{0.22em}
    \item Tiga kelompok bus: alamat 20-bit (satu arah), data 16-bit
          (dua arah), kontrol. Siklus instruksi:
          \eng{fetch}--\eng{decode}--\eng{execute}--\eng{store}.
\end{itemize}
\end{frame}

% ---------------------------------------------------------
\begin{frame}{Cek Pemahaman}
\small
\begin{enumerate}
    \item Unit yang menghitung alamat fisik 20-bit dari pasangan
          segmen:offset adalah \dots
          \begin{enumerate}\footnotesize
              \item \eng{Execution Unit}
              \item \eng{Bus Interface Unit}
              \item \eng{Arithmetic Logic Unit}
              \item register \eng{flag}
          \end{enumerate}

    \vspace{0.2em}
    \item Ukuran \eng{instruction queue} pada 8086 adalah \dots
          \begin{enumerate}\footnotesize
              \item 2 byte \quad \item 4 byte \quad
              \item 6 byte \quad \item 16 byte
          \end{enumerate}

    \vspace{0.2em}
    \item Bila \reg{CS} $=$ \hex{1000} dan \reg{IP} $=$ \hex{0025},
          instruksi berikutnya diambil dari alamat fisik \dots
          \begin{enumerate}\footnotesize
              \item \hex{10025} \quad \item \hex{10250} \quad
              \item \hex{1002} \quad \item \hex{100025}
          \end{enumerate}
\end{enumerate}
\end{frame}

% =========================================================
% PENUTUP
% =========================================================
\begin{frame}{Penutup}
\small
\begin{block}{Sebelum pertemuan berikutnya}
\begin{itemize}
    \item Selesaikan praktikum: diagram blok BIU/EU dari ingatan,
          tabel penelusuran, dan perbandingan jumlah siklus.
    \item Pelajari kembali hubungan panjang instruksi dengan
          kenaikan \reg{IP}.
    \item Baca \textbf{Bab V} --- Register 8086 dan Register Flag.
\end{itemize}
\end{block}

\vspace{0.5em}
\begin{center}
{\color{obeBlue}\sffamily\large\bfseries Terima kasih\par}
\vspace{0.2em}
{\color{obeGray}\sffamily\small \NamaMK\ (\KodeMK) \textbullet\
\SemesterTeks\par}
\end{center}
\end{frame}

\end{document}
